\documentclass[10pt,a4paper]{amsart}
\usepackage{amsmath,amssymb,amsthm,graphicx}
\usepackage{verbatim}
\newenvironment{reference}{\comment}{\endcomment}
%\newenvironment{reference}{}{}
\newenvironment{slogan}{\comment}{\endcomment}
\newenvironment{history}{\comment}{\endcomment}

% For commutative diagrams we use Xy-pic
\usepackage[all]{xy}

% We use 2cell for 2-commutative diagrams.
\xyoption{2cell}
% Theorem environments.
%
\theoremstyle{plain}
\newtheorem{theorem}[subsection]{Theorem}
\newtheorem{proposition}[subsection]{Proposition}
\newtheorem{lemma}[subsection]{Lemma}

\theoremstyle{definition}
\newtheorem{definition}[subsection]{Definition}
\newtheorem{example}[subsection]{Example}
\newtheorem{exercise}[subsection]{Exercise}
\newtheorem{situation}[subsection]{Situation}

\theoremstyle{remark}
\newtheorem{remark}[subsection]{Remark}
\newtheorem{remarks}[subsection]{Remarks}

\numberwithin{equation}{subsection}

% Macros
%
\def\lim{\mathop{\mathrm{lim}}\nolimits}
\def\colim{\mathop{\mathrm{colim}}\nolimits}
\def\Spec{\mathop{\mathrm{Spec}}}
\def\Hom{\mathop{\mathrm{Hom}}\nolimits}
\def\Ext{\mathop{\mathrm{Ext}}\nolimits}
\def\SheafHom{\mathop{\mathcal{H}\!\mathit{om}}\nolimits}
\def\SheafExt{\mathop{\mathcal{E}\!\mathit{xt}}\nolimits}
\def\Sch{\mathit{Sch}}
\def\Mor{\mathop{\mathrm{Mor}}\nolimits}
\def\Ob{\mathop{\mathrm{Ob}}\nolimits}
\def\Sh{\mathop{\mathit{Sh}}\nolimits}
\def\NL{\mathop{N\!L}\nolimits}
\def\CH{\mathop{\mathrm{CH}}\nolimits}
\def\proetale{{pro\text{-}\acute{e}tale}}
\def\etale{{\acute{e}tale}}
\def\QCoh{\mathit{QCoh}}
\def\Ker{\mathop{\mathrm{Ker}}}
\def\Im{\mathop{\mathrm{Im}}}
\def\Coker{\mathop{\mathrm{Coker}}}
\def\Coim{\mathop{\mathrm{Coim}}}

% Boxtimes
%
\DeclareMathSymbol{\boxtimes}{\mathbin}{AMSa}{"02}

%
% Macros for moduli stacks/spaces
%
\def\QCohstack{\mathcal{QC}\!\mathit{oh}}
\def\Cohstack{\mathcal{C}\!\mathit{oh}}
\def\Spacesstack{\mathcal{S}\!\mathit{paces}}
\def\Quotfunctor{\mathrm{Quot}}
\def\Hilbfunctor{\mathrm{Hilb}}
\def\Curvesstack{\mathcal{C}\!\mathit{urves}}
\def\Polarizedstack{\mathcal{P}\!\mathit{olarized}}
\def\Complexesstack{\mathcal{C}\!\mathit{omplexes}}
% \Pic is the operator that assigns to X its picard group, usage \Pic(X)
% \Picardstack_{X/B} denotes the Picard stack of X over B
% \Picardfunctor_{X/B} denotes the Picard functor of X over B
\def\Pic{\mathop{\mathrm{Pic}}\nolimits}
\def\Picardstack{\mathcal{P}\!\mathit{ic}}
\def\Picardfunctor{\mathrm{Pic}}
\def\Deformationcategory{\mathcal{D}\!\mathit{ef}}
\usepackage{cleveref}
\crefname{originalreference}{original reference}{original references}
\Crefname{originalreference}{Original reference}{Original references}
\makeatletter
\newlabel{categories-definition-exact}{{0034}{source}{}{sourcepointer-1}{}}
\newlabel{categories-definition-exact@cref}{{[originalreference][0][]0034}{[1][0][]source}{}{}{}}
\makeatother
\crefname{originallemma}{original lemma}{original lemmas}
\Crefname{originallemma}{Original lemma}{Original lemmas}
\makeatletter
\newlabel{lemma-functoriality-etale-site}{{03G2}{source}{}{sourcepointer-2}{}}
\newlabel{lemma-functoriality-etale-site@cref}{{[originallemma][0][]03G2}{[1][0][]source}{}{}{}}
\makeatother
\crefname{originalreference}{original reference}{original references}
\Crefname{originalreference}{Original reference}{Original references}
\makeatletter
\newlabel{sites-definition-morphism-sites}{{00X1}{source}{}{sourcepointer-3}{}}
\newlabel{sites-definition-morphism-sites@cref}{{[originalreference][0][]00X1}{[1][0][]source}{}{}{}}
\makeatother
\crefname{originalreference}{original reference}{original references}
\Crefname{originalreference}{Original reference}{Original references}
\makeatletter
\newlabel{sites-modules-lemma-flat-pullback-exact}{{04JC}{source}{}{sourcepointer-4}{}}
\newlabel{sites-modules-lemma-flat-pullback-exact@cref}{{[originalreference][0][]04JC}{[1][0][]source}{}{}{}}
\makeatother
\crefname{originalreference}{original reference}{original references}
\Crefname{originalreference}{Original reference}{Original references}
\makeatletter
\newlabel{sites-lemma-pullback-representable-sheaf}{{04D3}{source}{}{sourcepointer-5}{}}
\newlabel{sites-lemma-pullback-representable-sheaf@cref}{{[originalreference][0][]04D3}{[1][0][]source}{}{}{}}
\makeatother
\crefname{originalreference}{original reference}{original references}
\Crefname{originalreference}{Original reference}{Original references}
\makeatletter
\newlabel{sites-lemma-point-morphism-sites}{{05V1}{source}{}{sourcepointer-6}{}}
\newlabel{sites-lemma-point-morphism-sites@cref}{{[originalreference][0][]05V1}{[1][0][]source}{}{}{}}
\makeatother
\crefname{originallemma}{original lemma}{original lemmas}
\Crefname{originallemma}{Original lemma}{Original lemmas}
\makeatletter
\newlabel{lemma-stalk-gives-point}{{04K0}{source}{}{sourcepointer-7}{}}
\newlabel{lemma-stalk-gives-point@cref}{{[originallemma][0][]04K0}{[1][0][]source}{}{}{}}
\makeatother
\crefname{originallemma}{original lemma}{original lemmas}
\Crefname{originallemma}{Original lemma}{Original lemmas}
\makeatletter
\newlabel{lemma-stalk-exact}{{04K1}{source}{}{sourcepointer-8}{}}
\newlabel{lemma-stalk-exact@cref}{{[originallemma][0][]04K1}{[1][0][]source}{}{}{}}
\makeatother
\crefname{originalreference}{original reference}{original references}
\Crefname{originalreference}{Original reference}{Original references}
\makeatletter
\newlabel{sites-section-continuous-functors}{{00WU}{source}{}{sourcepointer-9}{}}
\newlabel{sites-section-continuous-functors@cref}{{[originalreference][0][]00WU}{[1][0][]source}{}{}{}}
\makeatother
\crefname{originalreference}{original reference}{original references}
\Crefname{originalreference}{Original reference}{Original references}
\makeatletter
\newlabel{sites-section-functoriality-PSh}{{00VC}{source}{}{sourcepointer-10}{}}
\newlabel{sites-section-functoriality-PSh@cref}{{[originalreference][0][]00VC}{[1][0][]source}{}{}{}}
\makeatother
\crefname{originalequation}{original equation}{original equations}
\Crefname{originalequation}{Original equation}{Original equations}
\makeatletter
\newlabel{equation-stalk-spaces-etale}{{04JZ}{source}{}{sourcepointer-11}{}}
\newlabel{equation-stalk-spaces-etale@cref}{{[originalequation][0][]04JZ}{[1][0][]source}{}{}{}}
\makeatother
\title{Pullbacks of sheaves and stalks along morphisms of algebraic spaces}
\author{The Stacks Project Authors}
\date{stacks-project:spaces-properties.tex:local-snapshot-20261012; GNU FDL 1.2 or later}
\begin{document}
\maketitle
\noindent\textit{Editorial note.} All delivered mathematical text is the author's own, copied verbatim from the pinned Stacks Project chapter file spaces-properties.tex; no statement or proof step was rewritten, reordered or removed.\par\medskip
\noindent\textit{Editorial note.} Span recomputation by label pairing (LANE\_BRIEF 10.2): in this chapter the author's proof environment follows the statement environment as a standalone block; the spec covers the complete environment pair (statement span, proof span) as recorded in the pool candidate.\par\medskip
\noindent\textit{Editorial note.} Occupation scan: combined claim\_key + statement\_anchor + generic-line scan over all registration files (active and withdrawn) found 13 S10-geom published cards for spaces-properties.tex (prop-occupied.json, 2026-10-14); those 13 labels are excluded from this batch by design. For the 77 free labels: 0 occupied. The registry mentions field (7557 files) and focused-seeds (645) both point to the same S10-geom cards (dual-source check), no other occupation.\par\medskip
\ifdefined\hypertarget\hypertarget{sourcepointer-1}{}\fi
\noindent\textit{Source reference.} Original reference 0034: sources/stacks-chow/tags.txt, source lines 121--121. Cross-chapter standard prerequisite cited by the author inside the counted span; it is used as a pointer and its own proof is not part of this item.\par\medskip
\ifdefined\hypertarget\hypertarget{sourcepointer-2}{}\fi
\noindent\textit{Source reference.} Original lemma 03G2: sources/stacks-chow/tags.txt, source lines 4513--4513. Cross-chapter standard prerequisite cited by the author inside the counted span; it is used as a pointer and its own proof is not part of this item.\par\medskip
\ifdefined\hypertarget\hypertarget{sourcepointer-3}{}\fi
\noindent\textit{Source reference.} Original reference 00X1: sources/stacks-chow/tags.txt, source lines 1203--1203. Cross-chapter standard prerequisite cited by the author inside the counted span; it is used as a pointer and its own proof is not part of this item.\par\medskip
\ifdefined\hypertarget\hypertarget{sourcepointer-4}{}\fi
\noindent\textit{Source reference.} Original reference 04JC: sources/stacks-chow/tags.txt, source lines 5933--5933. Cross-chapter standard prerequisite cited by the author inside the counted span; it is used as a pointer and its own proof is not part of this item.\par\medskip
\ifdefined\hypertarget\hypertarget{sourcepointer-5}{}\fi
\noindent\textit{Source reference.} Original reference 04D3: sources/stacks-chow/tags.txt, source lines 5713--5713. Cross-chapter standard prerequisite cited by the author inside the counted span; it is used as a pointer and its own proof is not part of this item.\par\medskip
\ifdefined\hypertarget\hypertarget{sourcepointer-6}{}\fi
\noindent\textit{Source reference.} Original reference 05V1: sources/stacks-chow/tags.txt, source lines 7535--7535. Cross-chapter standard prerequisite cited by the author inside the counted span; it is used as a pointer and its own proof is not part of this item.\par\medskip
\ifdefined\hypertarget\hypertarget{sourcepointer-7}{}\fi
\noindent\textit{Source reference.} Original lemma 04K0: sources/stacks-chow/tags.txt, source lines 5956--5956. Cross-chapter standard prerequisite cited by the author inside the counted span; it is used as a pointer and its own proof is not part of this item.\par\medskip
\ifdefined\hypertarget\hypertarget{sourcepointer-8}{}\fi
\noindent\textit{Source reference.} Original lemma 04K1: sources/stacks-chow/tags.txt, source lines 5957--5957. Cross-chapter standard prerequisite cited by the author inside the counted span; it is used as a pointer and its own proof is not part of this item.\par\medskip
\ifdefined\hypertarget\hypertarget{sourcepointer-9}{}\fi
\noindent\textit{Source reference.} Original reference 00WU: sources/stacks-chow/tags.txt, source lines 1196--1196. Cross-chapter standard prerequisite cited by the author inside the counted span; it is used as a pointer and its own proof is not part of this item.\par\medskip
\ifdefined\hypertarget\hypertarget{sourcepointer-10}{}\fi
\noindent\textit{Source reference.} Original reference 00VC: sources/stacks-chow/tags.txt, source lines 1142--1142. Cross-chapter standard prerequisite cited by the author inside the counted span; it is used as a pointer and its own proof is not part of this item.\par\medskip
\ifdefined\hypertarget\hypertarget{sourcepointer-11}{}\fi
\noindent\textit{Source reference.} Original equation 04JZ: sources/stacks-chow/tags.txt, source lines 5955--5955. Cross-chapter standard prerequisite cited by the author inside the counted span; it is used as a pointer and its own proof is not part of this item.\par\medskip
\begin{lemma}
\label{lemma-stalk-pullback}
Let $S$ be a scheme.
Let $f : X \to Y$ be a morphism of algebraic spaces over $S$.
\begin{enumerate}
\item The functor
$f_{small}^{-1} :
\textit{Ab}(Y_\etale)
\to
\textit{Ab}(X_\etale)$
is exact.
\item The functor
$f_{small}^{-1} :
\Sh(Y_\etale)
\to
\Sh(X_\etale)$
is exact, i.e., it commutes with finite limits and colimits, see
Categories, Definition \ref{categories-definition-exact}.
\item For any \'etale morphism $V \to Y$ of algebraic spaces
we have $f_{small}^{-1}h_V = h_{X \times_Y V}$.
\item Let $\overline{x} \to X$ be a geometric point.
Let $\mathcal{G}$ be a sheaf on $Y_\etale$.
Then there is a canonical identification
$$
(f_{small}^{-1}\mathcal{G})_{\overline{x}} = \mathcal{G}_{\overline{y}}.
$$
where $\overline{y} = f \circ \overline{x}$.
\end{enumerate}
\end{lemma}

\begin{proof}
Recall that $f_{small}$ is defined via $f_{spaces, small}$ in
Lemma \ref{lemma-functoriality-etale-site}.
Parts (1), (2) and (3) are general consequences of the fact that
$f_{spaces, \etale} :
X_{spaces, \etale}
\to
Y_{spaces, \etale}$
is a morphism of sites, see
Sites, Definition \ref{sites-definition-morphism-sites}
for (2),
Modules on Sites, Lemma \ref{sites-modules-lemma-flat-pullback-exact}
for (1), and
Sites, Lemma \ref{sites-lemma-pullback-representable-sheaf}
for (3).

\medskip\noindent
Proof of (4). This statement is a special case of
Sites, Lemma \ref{sites-lemma-point-morphism-sites}
via
Lemma \ref{lemma-stalk-gives-point}.
We also provide a direct proof. Note that by
Lemma \ref{lemma-stalk-exact}.
taking stalks commutes with sheafification.
Let $\mathcal{G}'$ be the sheaf on $Y_{spaces, \etale}$ whose restriction
to $Y_\etale$ is $\mathcal{G}$.
Recall that $f_{spaces, \etale}^{-1}\mathcal{G}'$ is the sheaf
associated to the presheaf
$$
U \longrightarrow \colim_{U \to X \times_Y V} \mathcal{G}'(V),
$$
see
Sites, Sections \ref{sites-section-continuous-functors} and
\ref{sites-section-functoriality-PSh}.
Thus we have
\begin{align*}
(f_{spaces, \etale}^{-1}\mathcal{G}')_{\overline{x}}
& =
\colim_{(U, \overline{u})} f_{spaces, \etale}^{-1}\mathcal{G}'(U)
\\
& = \colim_{(U, \overline{u})}
\colim_{a : U \to X \times_Y V} \mathcal{G}'(V) \\
& = \colim_{(V, \overline{v})} \mathcal{G}'(V) \\
& = \mathcal{G}'_{\overline{y}}
\end{align*}
in the third equality the pair $(U, \overline{u})$ and the map
$a : U \to X \times_Y V$ corresponds to the pair $(V, a \circ \overline{u})$.
Since the stalk of $\mathcal{G}'$
(resp.\ $f_{spaces, \etale}^{-1}\mathcal{G}'$)
agrees with the stalk of $\mathcal{G}$ (resp.\ $f_{small}^{-1}\mathcal{G}$),
see
Equation (\ref{equation-stalk-spaces-etale})
the result follows.
\end{proof}

\end{document}
